%% kapa: exam=1 round=26 year=2015 period=1 form=A subject=economics subject_ko=경제원론 seq=29 number=69
%% source: Q-Net 감정평가사 기출문제 (cst00309) artlSeq=5138001 file=extracted/1008/5138001/제26회_감정평가사 1차_기출/제26회_감정평가사 1차_기출/1교시_감평_A형.pdf
%% conversion: kordoc 4.24.51 -> tools/gen1.py (auto); review=proofread-sonnet; flags=none
\begin{question}{69}
A국 경제가 유동성함정(LM곡선이 수평)에 빠졌을 경우, 이에 관한 설명으로 옳은 것은?
\begin{choices}
  \item 투자가 이자율에 대해 매우 탄력적이다.
  \item 확대통화정책이 확대재정정책보다 국민소득을 더 많이 증가시킨다.
  \item 확대재정정책을 시행하면 구축효과로 인해 국민소득의 변화가 없다.
  \item 화폐수요가 이자율에 대해 완전비탄력적이다.
  \item 확대통화정책을 시행하더라도 이자율의 변화가 없다.
\end{choices}
\end{question}
\answer{5}
